\documentclass[../Tesi.tex]{subfiles}


\begin{document}
\section{Introduzione}
	
	\subsection{Scopo dello stage}
		L'\textbf{Experience API}(xAPI anche conosciuta come Tin Can API) è una specifica software per la creazione di contenuti per l'e-learning. È considerata il successore dello \textbf{SCORM}, ovvero lo standard de facto utilizzato per l'inserimento dei contenuti di e-learning all'interno di \textbf{Learning Management Systems}. Tali piattaforme sono necessarie per la distribuzione dei contenuti in formato SCORM e sono normalmente accessibili utilizzando un browser, per la visualizzazione dei contenuti in esse pubblicati. Con la nuova specifica, però, non sono più necessarie: i contenuti xAPI sono fruibili anche senza l'utilizzo di browser e di connessioni internet. Le informazioni sulle esperienze degli utenti sono tracciate utilizzando degli statement in formato JavaScript Object Notation(JSON) e sono raccolti da un \textbf{Learning Record Store}(LRS). Gli statement sono trasmessi agli LRS, utilizzando connessioni HTTP o HTTPS, che provvedono alla loro memorizzazione. Nella forma più semplice gli statement sono composti da tre parti: 
		\begin{itemize}
			\item \textit{attore}: specifica chi ha compiuto una determinata azione;
			\item \textit{verbo}: l'azione compiuta dall'attore;
			\item \textit{oggetto}: l'oggetto a cui è rivolta l'azione.
		\end{itemize}
		Sfruttando tale specifica si vuole creare un'applicazione Android che permetta la fruizione di contenuti per l'e-learning sia quando il dispositivo è online che offline. Quando il dispositivo non è provisto di una connessione internet attiva le esperienze di un utente dovranno essere salvate sul dispositivo in modo tale che, quando sarà nuovamente disponibile un accesso ad Internet, sia possibile inviare tali statement ad un LRS. L'applicazione deve inoltre presentare la possibilità di specificare quali contenuti devono essere visualizzati.

	\subsection{Azienda ospitante}
	ModoNetwork è una piccola azienda con sede a Caerano San Marco, nella provincia di Treviso, che opera nel campo del retail. Nasce nel 1998 per opera di Germana Boghetto, Luca Tonello e Andrea Comirato
	\begin{figure}[H]
		\centering
		\includegraphics[scale=0.1]{images/logoMN}
		\caption{Logo azienda ModoNetwork}
	\end{figure}
	L'azienda è particolarmente attiva nel campo della formazione. Essa si occupa, infatti, di seguire progetti di e-learning, a partire dalla creazione dei corsi fino, eventualmente, alla loro gestione all'interno di piattaforme LMS, prestando particolare attenzione alle esigenze del cliente e alla qualità del prodotto finale.

	\subsection{Struttura del documento}
		In questa parte viene riportata la struttura del documento. Ad ogni sezione è associata una breve descrizione del contenuto.
		\begin{enumerate}
			\item Introduzione: scopo dello stage e contenuti presenti nel documento;
			\item Progetto dello stage: descrizione degli obiettivi dello stage e delle caratteristiche che il prodotto finale deve avere;
			\item Tecnologie utilizzate: descrizione delle tecnologie utilizzate durante lo stage;
			\item Progettazione: descrizione delle scelte progettuali fatte riguardanti l'applicazione;
			\item Descrizione dei package: descrizione di tutti i package in cui è organizzato il codice del prodotto finale;
			\item Descrizione delle classe principali: descrizione delle classi più importanti dell'applicazione;
			\item Verifica e validazione: descrizione della fase di testing dell'applicazione;
			\item Conclusioni: bilancio finale sull'attività di stage.
		\end{enumerate}

\end{document}